\documentclass[12pt, a4paper]{article}

\usepackage[margin=2.5cm]{geometry} 
\usepackage{fancyhdr}               
\usepackage{hyperref}               
\usepackage{graphicx}               
\usepackage{float}                  
\usepackage{titlesec}
\usepackage{wrapfig}         
\usepackage{subcaption}

% Pengaturan Header dan Footer
\pagestyle{fancy}
\fancyhf{}
\lhead{\textbf{Nama:} Khalilullah Al Faath}   
\rhead{\textbf{NIM:} 25/566643/PPA/07139}        
\renewcommand{\headrulewidth}{0.4pt}         
\renewcommand{\footrulewidth}{0.4pt}         


\begin{document}

\begin{center}
    \large \textbf{Judul Laporan} \\[0.2cm]
    \large Nama Mata Kuliah  \\[0.5cm]
\end{center}


\section{Pranala Kode}
Hasil \textit{running} implementasi kode dapat diakses di pranala repositori GitHub berikut \url{https://github.com/khalilullahalfaath/comparing-ranked-retrieval-models/blob/main/1_STBI.ipynb}.


\end{document}